% =============================================================================
\section{Dimension 1: Workload Characterization}
\label{sec:workload}
% =============================================================================

Our fleet optimization work~\cite{poolrouting2026, fleetopt2026,
fleetsim2026} showed that workload characterization drives the remaining
design decisions.  We organize workloads
along four axes.

% -----------------------------------------------------------------------------
\subsection{Chat vs.\ Agent Workloads}
\label{sec:workload:chat-agent}
% -----------------------------------------------------------------------------

Chat workloads (LMSYS-Chat-1M~\cite{zheng2024lmsys},
the Azure LLM Inference Trace published alongside
Splitwise~\cite{patel2024splitwise}) are dominated by short prompts:
80--95\% of requests fall below 8K tokens, with a heavy tail to 64K+.
Agent workloads (SWE-bench~\cite{jimenez2024swebench},
BFCL~\cite{patil2025bfcl}) exhibit prompt lengths that grow
deterministically with turn number as tool outputs accumulate.
ServeGen~\cite{servegen2025} reports that agents make 3--10$\times$
more LLM calls per user request than chatbots.

Both types are \emph{multi-turn}, but with different session
structures.  Chat sessions are conversational: each turn
carries implicit context from prior exchanges, and the user may
shift topics, provide corrections, or gradually escalate adversarial
probes across turns.  Agent sessions are task-driven: context grows
monotonically as tool outputs accumulate, and the session terminates
on task completion or failure.  Multi-turn APIs (such as OpenAI's
Responses API) encode these patterns explicitly, offering
conversation-state management and turn-level optimizations that
single-request routing cannot exploit.  A multi-turn-aware router must maintain session context for two
reasons: KV-cache affinity and routing signal quality.  Prior turns
reveal user intent, topic drift, difficulty escalation, and safety
violations that a single-turn classifier would miss.

Our pool routing paper~\cite{poolrouting2026} first exposed how
the chat/agent distinction creates routing opportunities: the
80--95\% of short chat requests are over-provisioned in a
homogeneous 64K fleet, wasting $7/8$ of KV-cache slots.  The
deterministic turn-by-turn growth of agent sessions suggests that
session-aware prediction could outperform per-category EMA for
routing.

\paragraph{Memory management as a hidden cost multiplier.}
Beyond raw context growth, multi-turn sessions face a compounding
cost problem rooted in memory management failures.  When an LLM
hallucinates or selects the wrong tool, the erroneous output enters
the conversation history and is re-sent on every subsequent
turn---the organization pays input-token cost for content that
\emph{reduced} accuracy.  The agent must either tolerate the
polluted context (risking further errors through
self-conditioning~\cite{driftbench2026}) or spend additional LLM
calls to summarize, compress, or remove the bad turns before
continuing.  AgentHallu~\cite{agenthallu2026} benchmarks this
challenge: even frontier models achieve only 41.1\% accuracy at
localizing which step caused a hallucination, dropping to 11.6\%
for tool-use hallucinations---making targeted cleanup unreliable and
often requiring full-context re-summarization.

A comprehensive memory
survey~\cite{memorysurv2026} identifies five mechanism families for
agent memory (context-resident compression, retrieval-augmented
stores, reflective self-improvement, hierarchical virtual context,
and policy-learned management), each with distinct cost profiles.
A direct cost-performance comparison~\cite{memcost2026} shows that
at 100K-token context lengths, per-turn inference charges grow
proportionally with context even under prompt caching, and
fact-based memory systems become cheaper after approximately ten
interaction turns---precisely the regime where agent sessions
operate.  Context compression can mitigate the
problem: ACON~\cite{acon2025} achieves 26--54\% peak-token reduction
while preserving 95+\% task accuracy via failure-driven guideline
optimization, and Focus~\cite{focus2026} demonstrates 22.7\% token
savings through autonomous agent-driven context pruning on SWE-bench
tasks.  However, both operate on the agent side, optimizing a
single session in isolation.  This is a fundamental limitation:
SAMULE~\cite{samule2025} shows that \emph{single-trajectory}
reflection (the only level available to an individual agent)
produces narrow error corrections, while \emph{inter-task} learning
that aggregates failure patterns across diverse sessions yields
transferable insights that significantly outperform per-session
baselines.  CORRECT~\cite{correct2025} reinforces this: multi-agent
errors recur with similar structural patterns across requests, and
an online cache of distilled error schemata---built from prior
sessions---improves step-level error localization by up to 19.8\%.
Mistake Notebook Learning~\cite{mnl2025} demonstrates the same
principle: batch-clustered failure analysis produces generalizable
guidance that prevents known pitfalls without parameter updates.

The router has a structural advantage that directly exploits this
cross-session transferability: it observes outcomes from thousands
of sessions per domain across all tenants, aggregating failure
patterns into the per-\emph{(model, domain, failure-class)} table
(\Cref{opp:failure-routing}).  An individual agent, by contrast, sees
only its own session---too few examples to build reliable failure
statistics or discover domain-specific compression policies (e.g.,
``coding sessions tolerate aggressive early-turn pruning of failed
tool traces; medical sessions must retain full context for compliance
audits'').  The router provides this cross-session intelligence as a
fleet-wide service that every agent benefits from immediately,
including new agents that would otherwise face a cold-start problem.

\paragraph{Gateway-native assistants and fleet observability.}
OpenClaw~\cite{openclaw2026} is a concrete open-source example of the
pattern \emph{personal assistant + local-first Gateway}: users talk
to the agent over many messaging channels while a WebSocket Gateway
coordinates sessions, tool execution (browser, automation, device
nodes), and model calls.  Such deployments generate long, messy
multi-turn traces---tool outcomes, retries, channel switches, and
compaction events---that fixed academic benchmarks rarely capture.
The inference router sits on the hot path for model traffic from
those gateways and can log which model, tool surface, and safety path
correlated with short vs.\ runaway sessions.  Aggregating traces
across tenants supplies offline RL and fleet-prior datasets that a
single local assistant cannot
amass~(\Cref{opp:session-cascade}, \Cref{opp:tool-pruning}).
Conversely, routing policies tuned for token- and round-minimization
feed back into operators' choices of model defaults and tool budgets
at the Gateway.

\paragraph{Proxy-level instruction, tool, and memory surfaces.}
ITR~\cite{itr2026} retrieves minimal system-instruction fragments
and a narrowed tool subset per agent step; on the authors' controlled
agent benchmark it reports $\sim$95\% lower per-step context tokens
and large gains in correct tool routing because monolithic catalogs
can consume $\sim$90\% of the context window.  The same mechanism can be implemented at an
inference \emph{gateway}---rewriting \texttt{tools} payloads and
system blocks before the request reaches the model---so frameworks
that always send full catalogs still benefit from fleet-wide
calibration.  ToolScope~\cite{toolscope2025} merges and filters tools
on the agent side; SMART~\cite{smart2026} shows that curbing tool
overuse improves both cost and downstream accuracy.  A router that
coordinates narrowing, merging hints, and budget signals
(BATS~\cite{bats2025}) avoids duplicating incompatible policies in
every SDK.  Finally, memory is not only ``what the agent
remembers'' but \emph{what the next prefill pays for}: MemCost-style
analyses~\cite{memcost2026} show when fact stores beat raw context,
and session-aware schedulers (AgServe~\cite{agserve2025},
Continuum~\cite{continuum2025}) treat KV-cache lifetime and
multi-turn structure as first-class resources.  The router is the
natural place to choose among full history, injected summaries, and
retrieval-first prefill for the next turn---reducing tokens
\emph{and} the error compounding that lengthens
rounds~(\Cref{opp:gateway-agent-loops}).

% -----------------------------------------------------------------------------
\subsection{Warm vs.\ Cold Requests}
\label{sec:workload:warm-cold}
% -----------------------------------------------------------------------------

\emph{Warm} requests hit a KV-cache prefix from a prior turn or
shared system prompt; \emph{cold} requests start from scratch.
Cached tokens cost up to $10\times$ less~\cite{llmd2026kvcache},
and prefix-cache-aware scheduling achieves $57\times$ faster
response times in distributed deployments~\cite{llmd2026kvcache}.

For agentic workloads, 40--60\% of session time involves paused tool
calls, causing LRU eviction to collapse cache hit
rates~\cite{vllm2026rfc37003}.  Session-affinity routing---keeping a
session on the same pool instance---could preserve KV-cache locality.
Category-aware semantic caching~\cite{wang2025cache} complements this
at a higher level: domain-aware similarity thresholds and TTLs
eliminate redundant inference requests entirely---returning cached
responses for semantically equivalent queries without invoking the
LLM, reducing both cost and load on the serving pool.


% -----------------------------------------------------------------------------
\subsection{Prefill-Heavy vs.\ Decode-Heavy Workloads}
\label{sec:workload:prefill-decode}
% -----------------------------------------------------------------------------

Vision-language models encode images into hundreds of visual tokens,
making VLM workloads \emph{prefill-heavy}.  Our AVR
system~\cite{avr2026} confronted this directly: CUA screenshots are
$\sim$5\,MB, and routing 4K screenshots requires multimodal signal
extraction.  FastRouter~\cite{fastrouter2026} showed that classical
NLP compression can cap router input to $\sim$512 tokens regardless
of original prompt length, making prefill-heavy workloads tractable
for the routing layer itself.

Coding agents are \emph{decode-heavy}: code generation produces long
structured outputs.  The decode ratio directly affects whether
disaggregated prefill/decode architectures
(Splitwise~\cite{patel2024splitwise},
DistServe~\cite{zhong2024distserve}) provide benefit, and our
inference-fleet-sim~\cite{fleetsim2026} tool models this trade-off
explicitly.


% -----------------------------------------------------------------------------
\subsection{Workload Identity Signals}
\label{sec:workload:identity}
% -----------------------------------------------------------------------------

Our prior work treats incoming requests as opaque text and infers
workload type through content analysis---embedding similarity,
keyword matching, and domain classifiers.  In practice, the request
structure itself carries rich identity signals that the router can
exploit without content inspection, opening the door to
zero-latency workload classification.

\paragraph{Explicit API-level signals.}
Modern inference APIs already expose structural metadata that
distinguishes agent from chat workloads:
\begin{itemize}[nosep]
  \item \textbf{Tool definitions.}  A request carrying a \texttt{tools}
        array with function schemas is almost certainly an agent
        workload; the number and types of tools (code execution,
        web search, MCP connectors) further classify the agent's
        operational mode.  The tool catalog itself is a cost lever:
        ITR~\cite{itr2026} shows that dynamically exposing only the
        minimal tool subset per step reduces per-step context tokens
        by 95\% and improves correct tool routing by 32\%, because
        monolithic tool catalogs consume up to 90\% of available
        context.
  \item \textbf{Session chaining.}  Fields such as
        \texttt{previous\_response\_id} (OpenAI Responses API) or
        conversation-state objects signal multi-turn continuations.
        Their presence tells the router the request is mid-session
        (enabling KV-cache retention, \Cref{opp:kv-retention}) and
        provides the turn number without parsing conversation history.
  \item \textbf{System-prompt fingerprints.}  System prompts are
        typically static per deployment and can be hashed at the
        gateway.  A lookup table of known hashes maps each deployment
        to its domain, expected tool set, and historical session-length
        distribution---giving the router a per-request prior at
        zero classification cost.
  \item \textbf{Gateway metadata headers.}  Production gateways
        (Cloudflare AI Gateway's \texttt{cf-aig-metadata},
        LiteLLM's \texttt{x-litellm-tags}, and custom
        \texttt{X-Client-Request-Id} headers) already let operators
        annotate requests with environment, team, and workload-type
        tags.  SIRP~\cite{sirp2025} standardizes this with
        classification-axis header fields.
\end{itemize}

\paragraph{Caller identity and authorization scope.}
The four signal classes above describe \emph{what} the request is.
None describes \emph{who} is making it or \emph{what they are
authorized to do}.  The bearer token carried in every LLM API call
fills this gap: in production gateways, it encodes caller identity
\emph{and} authorization scope---which models a key can access, which
tools it may invoke, and what budget it may consume.
\begin{itemize}[nosep]
  \item \textbf{Per-key tool permissions.}
        LiteLLM's Tool Permission Guardrail~\cite{litellm2026toolguard}
        enforces per-key tool allow/deny rules with regex matching on
        tool names and arguments.  Two modes---\emph{block} (reject the
        request) and \emph{rewrite} (strip forbidden tools before the
        model sees them)---demonstrate that caller-scoped tool governance
        is already production-ready at the gateway layer.
  \item \textbf{Structured agent authorization.}
        The Agent Authorization Profile
        (AAP)~\cite{aap2026}---an emerging OAuth~2.0 extension---adds
        five JWT claims: agent identity, capabilities with enforceable
        server-side constraints (domain restrictions, rate limits, time
        windows per tool), task binding (prevents scope drift), delegation
        tracking, and audit trails.  Its core principle is that
        ``authorization must be evaluated at execution time, not just
        during connection.''
  \item \textbf{Infrastructure analogue.}
        Kubernetes RBAC with admission webhooks enforces per-identity
        permissions at the API-server boundary regardless of what the
        application inside the pod believes it should do---the same
        zero-trust principle~\cite{nist800207} applied to LLM
        requests.
\end{itemize}
The bearer token enables a new class of routing decisions absent
from the current WRP roadmap: role-based tool-call enforcement
(\Cref{opp:rbac-enforcement}), per-caller fleet-wide reputation
scoring (\Cref{opp:caller-reputation}), and router-enforced
behavioral commitments (\Cref{opp:behavioral-commitments}).  It
also extends existing opportunities: per-caller cumulative risk
accumulation (\Cref{opp:cumulative-risk}), cross-session token
budgets (\Cref{opp:token-budget}), RBAC predicates in the
governance DSL (\Cref{opp:governance}), role-aware tool catalog
shaping (\Cref{opp:tool-pruning}), and caller-aware memory
isolation (\Cref{opp:gateway-agent-loops}).

\paragraph{Inferred structural signals.}
Beyond explicit metadata, the router can derive workload identity
from request structure at dispatch time:
\begin{itemize}[nosep]
  \item \textbf{Conversation depth.}  The number of message objects
        in the \texttt{messages} array reveals session maturity:
        a 2-message exchange is likely a cold-start chat; a
        20-message array with alternating \texttt{assistant}/\texttt{tool}
        roles is a deep agent session with different model, pool,
        and compression requirements.
  \item \textbf{Tool-call history pattern.}  If prior assistant
        messages contain \texttt{tool\_calls} with high failure rates
        (parseable from subsequent \texttt{tool} role messages
        reporting errors), the router can preemptively route to a
        stronger model or apply compression before the next turn.
  \item \textbf{Reasoning metadata.}  NofT~\cite{noft2026} shows
        that the number of chain-of-thought reasoning steps predicts
        task difficulty pre-inference, enabling difficulty-aware model
        selection with 95\% adversarial detection accuracy.
\end{itemize}

\paragraph{From static to per-stage routing.}
These signals enable a shift from per-request routing (one decision
per API call) to \emph{per-stage} routing within agentic workflows.
Aragog~\cite{aragog2025} demonstrates this: by decoupling a one-time
accuracy-preserving configuration step from a per-stage scheduler
that uses current system observations, it achieves 50--217\%
throughput improvement and 32--79\% latency reduction at peak load.
The WRP framework is well-positioned for this evolution: the
router already maintains session state for multi-turn awareness
(VCD~\cite{vcd2026}, Feedback Detector~\cite{feedbackdet2026}),
and the per-stage identity signals listed above provide the input
features for a stage-aware routing policy that adapts model selection,
compression, and pool assignment as the session evolves.


% -----------------------------------------------------------------------------
\subsection{Workload CDF Archetypes}
\label{sec:workload:archetypes}
% -----------------------------------------------------------------------------

FleetOpt~\cite{fleetopt2026} identified three prompt-length CDF
archetypes that determine optimal pool and routing configuration:

\begin{enumerate}[nosep]
  \item \textbf{Concentrated-below} (\eg Azure): 80--95\% short
        requests.  Aggressive compression minimizes long-pool cost.
  \item \textbf{Dispersed} (\eg LMSYS): requests span the length
        spectrum.  Balanced two-pool split with moderate compression.
  \item \textbf{Concentrated-above} (\eg SWE-bench): most requests
        are long.  Raise $\boundary_{\text{short}}$ rather than
        compress.
\end{enumerate}

These archetypes are not static: as organizations shift from chat to
agents, the CDF migrates from Archetype~1 toward Archetype~3.
FleetOpt~\cite{fleetopt2026} showed that the minimum-cost fleet is
qualitatively different for each archetype, and
inference-fleet-sim~\cite{fleetsim2026} confirmed that monolithic
topology can outperform disaggregated for Archetype~1 when
$\boundary_{\text{short}} < 4$K.


